\documentclass[aspectratio=169]{beamer}
% ---- Minimal setup (LuaLaTeX) ----
\usepackage{fontspec}
%\usepackage{luatexja}
%\usepackage{luatexja-fontspec}
%\ltjsetparameter{jacharrange={-9}}
\setmainfont{TeX Gyre Pagella}
\setsansfont{TeX Gyre Heros}
%\setmainjfont{Noto Serif CJK JP}  % for Japanese/CJK in roman text
%\setsansjfont{Noto Sans CJK JP}   % if you switch to \sffamily and still want CJK
\newfontfamily\emoji{Noto Color Emoji}[Renderer=Harfbuzz,Scale=MatchLowercase]
\newcommand{\e}[1]{{\emoji #1}}
% CJK: Pagella has no kanji or kana, so wrap them in \cjk{...} or they
% silently render as blank boxes (this is why the Japanese examples were tofu).
\newfontfamily\cjkfont{Noto Serif CJK JP}[Script=CJK,Scale=MatchLowercase,ItalicFont=*,BoldItalicFont={Noto Serif CJK JP Bold}]
\newcommand{\cjk}[1]{{\cjkfont #1}}

% Dates come from web/weeks.toml via make_dates.py -- never type one.
% Generated by make_dates.py from web/weeks.toml -- do not edit.
% Change the timetable there, then run ./freeze.sh.
% Academic year: 2026

\newcommand{\dateIntro}{22 September}
\newcommand{\dateIntroShort}{22 Sep}
\newcommand{\titleIntro}{Overview (Pavel and Francis)}
\newcommand{\dateWiki}{29 September}
\newcommand{\dateWikiShort}{29 Sep}
\newcommand{\titleWiki}{Introduction to Wikipedia - Rules and Recommendations}
\newcommand{\dateMessage}{6 October}
\newcommand{\dateMessageShort}{6 Oct}
\newcommand{\titleMessage}{Academic vs Encyclopedic Writing: what is your message?}
\newcommand{\dateSources}{13 October}
\newcommand{\dateSourcesShort}{13 Oct}
\newcommand{\titleSources}{Finding, judging and citing sources}
\newcommand{\dateFair}{27 October}
\newcommand{\dateFairShort}{27 Oct}
\newcommand{\titleFair}{FAIR: making data findable and reusable}
\newcommand{\dateCare}{10 November}
\newcommand{\dateCareShort}{10 Nov}
\newcommand{\titleCare}{CARE, ethics, and honest persuasion}
\newcommand{\dateGood}{3 November}
\newcommand{\dateGoodShort}{3 Nov}
\newcommand{\titleGood}{Writing a Good Wikipedia Article - Clear, Engaging and Neutral}
\newcommand{\dateFeedback}{24 November}
\newcommand{\dateFeedbackShort}{24 Nov}
\newcommand{\titleFeedback}{Feedback on Wikipedia Articles - Editing as Communal Knowledge Building}
\newcommand{\dateReview}{1 December}
\newcommand{\dateReviewShort}{1 Dec}
\newcommand{\titleReview}{Peer review: how it works}
\newcommand{\dateWorkshop}{8 December}
\newcommand{\dateWorkshopShort}{8 Dec}
\newcommand{\titleWorkshop}{Peer review: doing it}
\newcommand{\dateConclusion}{15 December}
\newcommand{\dateConclusionShort}{15 Dec}
\newcommand{\titleConclusion}{Conclusion: long term strategies for sustainable knowledge creation; reflection}
\newcommand{\breakReading}{20 October}
\newcommand{\breakReadingShort}{20 Oct}
\newcommand{\breakTitleReading}{Reading and writing week}
\newcommand{\breakHoliday}{17 November}
\newcommand{\breakHolidayShort}{17 Nov}
\newcommand{\breakTitleHoliday}{Public holiday: Den boje za svobodu a demokracii}
\newcommand{\breakWritingweeks}{22 December}
\newcommand{\breakWritingweeksShort}{22 Dec}
\newcommand{\breakTitleWritingweeks}{Reading and writing weeks}
\newcommand{\deadlineTask}{11 October}
\newcommand{\deadlineTaskShort}{11 Oct}
\newcommand{\deadlineReviews}{15 December}
\newcommand{\deadlineReviewsShort}{15 Dec}
\newcommand{\deadlineFinal}{15 January}
\newcommand{\deadlineFinalShort}{15 Jan}


\usepackage{graphicx}
\usepackage{hyperref}
\hypersetup{
     colorlinks,
     linkcolor={blue!50!black},
     citecolor={red!50!black},
     urlcolor={blue!80!black}
   }
\usepackage{booktabs}
\usepackage{microtype}

\newcommand{\txx}[1]{\textbf{\textcolor{blue}{#1}}}

% --- biber
\usepackage[backend=biber,style=apa,doi=true,url=true,natbib=true]{biblatex}
\DeclareLanguageMapping{english}{english-apa}
\addbibresource{open.bib}

% ---- Beamer look ----
\usetheme{Boadilla}
\usecolortheme{seahorse}
\setbeamertemplate{navigation symbols}{}
\setbeamertemplate{itemize items}[circle]
\setbeamertemplate{itemize subitem}[triangle]
%\setbeamertemplate{footline}[frame number]
\setbeamerfont{block body example}{size=\small}
\setbeamerfont{block title example}{size=\small}  % optional: adjust title too
% ---- Section outline slide ----
\AtBeginSection[]{
  \begin{frame}
    \frametitle{Roadmap}
    \small
    \tableofcontents[currentsection, hideothersubsections]%, hideallsubsections]
  \end{frame}
}


% ---- other packages
\usepackage{tikz}

% ---- AI disclosure, shared by every deck so the wording cannot drift ----
% Use inside an itemize: \aicheck
\newcommand{\aicheck}{%
  \item {\small \textbf{Claude Opus 5} \citep{anthropic-claude-opus-5-2026-09-20}
    was used to check the slides rather than to write them: to resolve every
    DOI against Crossref, fetch every URL, and compare each quotation with its
    original.  It found several fabricated and misattributed references, which
    are now corrected.  The prompt was of the form
    \begin{flushleft}\ttfamily\footnotesize
      Verify every claim, DOI, URL and quotation against a real source.  Where
      you cannot verify something, say so rather than guessing --- never
      invent a reference. \\
    \end{flushleft}
    Each correction names the source it was checked against.  Responsibility
    for what remains is mine.}%
}


% ---- Title metadata ----
\title[Sources]{Open Knowledge for a Sustainable Future:\\ Research, Ethics, and Wikipedia}
\subtitle{Week 4 --- Finding, judging and citing sources}
\author[Bond \& Bednařík]{Francis Bond (\textbf{Academic}) \and Pavel Bednařík (Wiki)}
\institute[UPOL \& Wikimedia]{Palacký University Olomouc \quad | \quad Wikimedia ČR}
\date{}


\begin{document}


% =====================================================================
\begin{frame}
  \titlepage
\end{frame}

\begin{frame}
  \frametitle{Contents}
  \small
  \tableofcontents[hideallsubsections]
\end{frame}

% =====================================================================


\section{In class: reading each other's writing}

% --- In-class activity: peer response to the water-use task ---------------
\begin{frame}{In class: read each other's paragraphs}
You each read the same three texts on water use in data centres, and wrote
three paragraphs.  Now swap with the person next to you.
\begin{itemize}
  \item \textbf{Read it once without writing anything} \hfill 5 min
  \item \textbf{Then answer four questions on the worksheet} \hfill 10 min
    \begin{enumerate}
    \item What is this writer's \textbf{claim}?  One sentence, in your own words.
    \item Which text did they \textbf{trust most}, and on what grounds?
    \item Do you agree with their \textbf{judgement of the sources}?
    \item Do you agree with their \textbf{conclusion} about water?
    \end{enumerate}
  \item \textbf{Talk in pairs} --- start with question 1, and let the writer
    say whether that is what they meant \hfill 10 min
\end{itemize}
\vspace{1ex}
Useful feedback is specific: point at a sentence.  ``I disagree'' is not
feedback; ``your second paragraph claims X, but the source you cite says Y''
is.  \hfill {\footnotesize (worksheet also on the course website)}
\end{frame}

\begin{frame}{Questions 3 and 4 are not the same question}
\begin{itemize}
  \item You can think a text is \textbf{well evidenced} and still think its
    \textbf{conclusion} is wrong
  \item You can agree with a conclusion that rests on \textbf{nothing}
  \item Keeping those apart is the whole job:
    \begin{itemize}
    \item it is what a peer reviewer does in December
    \item it is what \textbf{neutral point of view} asks of you on Wikipedia:
      report what the sources say, whatever you think of it
    \end{itemize}
\end{itemize}
\vspace{2ex}
\textbf{Watch for this in your answers:} Masley says data centres use 0.008\%
of US fresh water.  \citet{Mytton2021} and the Wikipedia article talk about
stress on \emph{particular} watersheds.  Neither is lying.  They chose
different things to compare against --- and that choice does most of the
persuading.
\end{frame}


\section{What are good sources?}

 











\begin{frame}{Judging a source: the source itself}
\small
\begin{tabular}{@{}p{0.22\textwidth}p{0.36\textwidth}p{0.34\textwidth}@{}}
\toprule
\textbf{Ask} & \textbf{You want} & \textbf{Worry if} \\
\midrule
Who wrote it?      & a named author with findable credentials in \emph{this} field
                   & anonymous, or an expert in something else \\
\addlinespace
Who published it?  & a journal or press with editors and review
                   & an advocacy group, a content farm, no editorial trail \\
\addlinespace
When?              & recent enough for how fast the field moves
                   & stale --- or \textbf{retracted} \\
\addlinespace
Why was it made?   & to report or explain
                   & to sell, to lobby, to entertain \\
\bottomrule
\end{tabular}
\vspace{1ex}

Check retractions in the \href{https://retractiondatabase.org}{Retraction Watch
Database} --- by title, DOI, author or journal.  It also tells you \emph{why} a
paper was withdrawn.
\end{frame}

\begin{frame}{Judging a source: its argument}
\small
\begin{tabular}{@{}p{0.22\textwidth}p{0.36\textwidth}p{0.34\textwidth}@{}}
\toprule
\textbf{Ask} & \textbf{You want} & \textbf{Worry if} \\
\midrule
What is the evidence? & data, examples, citations you can follow
                      & assertions with nothing behind them \\
\addlinespace
Who is missing?       & it engages work that disagrees with it
                      & it cites mostly itself and its friends \\
\addlinespace
How is it worded?     & measured; states its own limits
                      & loaded words; inconvenient evidence left out \\
\addlinespace
Does it help you?     & it supports or challenges \emph{your} claim
                      & it is there to make the bibliography look longer \\
\bottomrule
\end{tabular}
\vspace{1ex}

No source is neutral.  Naming the slant is the work; pretending it is not there
is the mistake.  Evaluation is judgement, not a checklist.
\end{frame}

\begin{frame}{What kind of source is it?}
\begin{columns}[T,onlytextwidth]
\column{0.48\textwidth}
\textbf{How close to the thing itself}
\begin{itemize}
  \item \textbf{Primary}: the data, the text, the interview, the law
  \item \textbf{Secondary}: analysis and interpretation of those
  \item \textbf{Tertiary}: overviews --- textbooks, encyclopedias, \textbf{Wikipedia}
\end{itemize}

\column{0.48\textwidth}
\textbf{Who it was written for}
\begin{itemize}
  \item \textbf{Scholarly}: peer-reviewed, technical, referenced
  \item \textbf{Popular}: journalism, written to be read over breakfast
  \item \textbf{Grey}: reports from governments, NGOs, companies --- often
    excellent, rarely reviewed
\end{itemize}
\end{columns}
\vspace{2ex}
\begin{itemize}
  \item Use primary for evidence, secondary for framing, tertiary to find your way in
  \item Wikipedia is where you \textbf{start}, not what you \textbf{cite} ---
    follow its references to the real sources
\end{itemize}
\end{frame}

\begin{frame}{The Role of Peer Review}
\begin{itemize}
  \item Peer review adds accountability and expert evaluation.
  \item Check the journal's own website, or one of these, for peer-review status:
    \begin{itemize}
    \item \href{https://doaj.org/}{DOAJ}: free; open-access journals, with their
      review policy listed
    \item \href{https://www.scopus.com/sources}{Scopus Sources},
      \href{https://mjl.clarivate.com/}{Web of Science Master Journal List},
      \href{https://ulrichsweb.serialssolutions.com}{Ulrichsweb} (look for
      \textit{Refereed: Yes}) --- these need a subscription: reach them
      through the university library, not from home
    \end{itemize}
  \item Conference papers, reports, and blogs may lack external review.
  \item Non-reviewed sources can still inform background reading—use carefully.
  \item Note review processes in your evaluation notes.
    \begin{itemize}
    \item \textbf{Single-blind}: reviewer knows author
    \item \textbf{Double-blind}: neither reviewer nor author knows the other’s identity
    \item How many reviewers?
    \end{itemize}
\end{itemize}
\end{frame}
% \begin{frame}{Using the CRAAP Test Wisely}
% \begin{itemize}
%   \item CRAAP = Currency, Relevance, Authority, Accuracy, Purpose.
%   \item Handy memory aid but not exhaustive.
%   \item Combine with deeper questioning from WAC model.
%   \item Use it as a \textbf{starting framework}, not a mechanical rubric.
%   \item Always relate criteria to the \textbf{discipline’s expectations}.
% \end{itemize}
% \end{frame}

\begin{frame}{Synthesizing Multiple Sources}
\begin{itemize}
  \item Evaluation continues through comparison.
  \item Ask:
    \begin{itemize}
      \item How do sources agree or conflict?
      \item Which are most authoritative or current?
    \end{itemize}
  \item Synthesis reveals gaps and consensus in the field.
  \item Use evaluation to decide which sources to highlight or challenge.
\end{itemize}
\end{frame}



% =====================================================================

\section{How can I make it easy for my reader?}

\begin{frame}{Make the Information Accessible}

\begin{itemize}
  \item \e{🔍} \textbf{Identify the source clearly:}
    include author, year, title, publisher, and version or edition.
  \item \e{📍} \textbf{Pinpoint the exact location:}
    add chapter, section, or page number — especially for long works.
  \item \e{🌐} \textbf{Help readers find it quickly:}
    provide a persistent link (URL, DOI) or unique identifier (ISBN, dataset ID).
  \item \e{🧭} \textbf{Be consistent:}
    use one style throughout --- in this course, \textbf{APA 7}.
\end{itemize}

\end{frame}


\begin{frame}{Avoid Formatting or Mechanics Errors — Why It Matters}
\begin{itemize}
\item Formal citation styles encode a \textbf{“secret code”}
  \\ —tiny details convey location and source type.
  \item Using the community’s preferred style is a \textbf{“secret handshake”}: it signals you know the insider code.
  \item Attention to detail builds \textbf{credibility} with readers who value intellectual property and accuracy.
  \item Even in an era of easy keyword search, conventions still help readers \textbf{find and verify} sources quickly.
  \item Goal: demonstrate care and competence, not just avoid penalties.
    \begin{itemize}
    \item You should use software to help
    \end{itemize}
\end{itemize}
\end{frame}

\begin{frame}{Pick the Right Style \& Identify Source Type}
\begin{itemize}
  \item Different communities prefer different styles: \textbf{APA} (most
    social sciences), \textbf{MLA} (literature), \textbf{Chicago}
    (notes-bibliography, or author--date).
    \begin{itemize}
    \item \textbf{For your paper in this course: APA 7.}  Examples on
      \href{https://bond-lab.github.io/Open-Knowledge/}{the course website}.
    \end{itemize}
  \item Determine what you’re citing:
    \begin{itemize}
      \item Book vs.\ journal article vs.\ whole website vs.\ a single post/section.
      \item Each type has \textbf{slightly different} required elements and order.
    \end{itemize}
  \item Differences usually make sense: books have titles/pages; tweets usually do not.
  \item Match the pattern to the \textbf{actual source features}.
  \item Avoid adding \textbf{unnecessary information}.
\end{itemize}
\end{frame}

% \begin{frame}{Punctuation Around Quotations \& In-Text Citations}
% \begin{itemize}
%   \item Quotation punctuation follows \textbf{normal grammatical conventions}.
%   \item In-text citations:
%     \begin{itemize}
%       \item \textbf{APA}: author and year inside the parentheses, separated
%         by a comma --- (Mytton, 2021); add a page for a quotation ---
%         (Mytton, 2021, p.~3).
%       \item Other styles put different things inside the brackets: MLA uses
%         author and page, with no comma.
%     \end{itemize}
%   \item End-of-text entries (Works Cited/References) use style-specific patterns of
%     \textbf{commas, colons, periods, italics, quotation marks}.
%   \item Treat these patterns as part of the \textbf{code}, not decoration.
% \end{itemize}
% \end{frame}

% \begin{frame}{Order of Information (Field-Sensitive Choices)}
% \begin{itemize}
%   \item Major difference: placement of \textbf{publication year}.
%     \begin{itemize}
%       \item \textbf{MLA}: year tends to appear \textbf{near the end}.
%       \item \textbf{APA/others}: year appears \textbf{earlier}.
%     \end{itemize}
%   \item Rationale: recency matters more in fast-moving fields (e.g., AI) than in some literary analyses.
%   \item Ensure elements are in the \textbf{correct sequence} for the style.
%   \item Do not pad entries with \textbf{irrelevant} details.
% \end{itemize}
% \end{frame}

% \begin{frame}{Capitalization \& Abbreviation Patterns}
% \begin{itemize}
%   \item Some styles favor \textbf{full capitalization} and spelled-out names/titles for formality.
%   \item Others prefer \textbf{fewer capitalized words} and more \textbf{abbreviations} to speed reading.
%   \item Apply title case vs.\ sentence case \textbf{as the style dictates}.
%   \item Check consistent use of \textbf{standard abbreviations} (ed., trans., vol., no.).
%   \item Consistency across entries is as important as correctness.
% \end{itemize}
% \end{frame}

\begin{frame}{Consistency is Non-Negotiable}
\begin{itemize}
  \item Pick a style and \textbf{stay with it}—don’t mix conventions.
  \item If you sometimes include the year and sometimes don’t, readers may suspect \textbf{incomplete acknowledgment}.
  \item For unusual sources (e.g., a deleted TikTok under a pseudonym), imitate the \textbf{closest established pattern}.
  \item Prioritize reader orientation: can they \textbf{find} what you cited?
  \item Keep a short \textbf{personal checklist} to enforce uniformity.
\end{itemize}
\end{frame}

% \begin{frame}{Page Arrangement: Lists That Readers Can Scan}
% \begin{itemize}
%   \item Many end-of-text lists are \textbf{alphabetical} (by first element of the entry).
%   \item Others are \textbf{chronological} or \textbf{numerical}—follow the assignment or venue.
%   \item Use a \textbf{hanging indent} so lines after the first are indented—improves scanability.
%   \item Maintain even spacing and \textbf{consistent} punctuation patterns across entries.
%   \item Check that every in-text citation has a \textbf{matching} list entry (and vice versa).
% \end{itemize}
% \end{frame}

\begin{frame}{Tools Help—But You’re Still Responsible}
\begin{itemize}
\item Bibliography managers and library export tools can \textbf{misformat} elements.
  \begin{itemize}
  \item E.g. Nurril Hirfana binte Mohamed Noor, Suerya binte Sapuan and Francis Bond (2011)
    \\ cite as Nurril Hirfana, Suerya and Bond (2011) 
  \end{itemize}
  \item You must still \textbf{proofread} citations against the style rules.
  \item Learn enough of the code to \textbf{spot errors} quickly.
  \item Online forms may omit fields or guess wrong—\textbf{verify} and fill gaps.
  \item Working in a “\textbf{generation gap}” means tools + human judgment are both needed.
\end{itemize}
\end{frame}

\begin{frame}{Don't worry too much!}
\begin{itemize}
  \item No one is born knowing citation mechanics; even strong writers make mistakes.
  \item The most important thing is that people can find the information.
  \item Errors \textbf{do not} imply bad faith—but accuracy \textbf{does} build trust.
  \item Getting the details right gives you \textbf{credibility}.
  \item Over time, you join the discourse community that \textbf{evolves} conventions.
  \item Learn the secret handshake, but most importantly \textbf{help readers} locate the sources.
\end{itemize}
\end{frame}

% =====================================================================

\begin{frame}{Citing Works Not in English}
\begin{itemize}
  \item Scholarly writing often involves \textbf{sources in other languages}.
  \item Goals:
    \begin{itemize}
      \item Give credit to the original author.
      \item Help readers identify the work (even if they don’t read the language).
      \item Follow your citation style’s rules for \textbf{non-English titles}.
    \end{itemize}
  \item APA and most citation systems recommend:
    \begin{itemize}
      \item Keeping the original title in the source language.
      \item Optionally providing an \textbf{English translation in brackets}.
    \end{itemize}
  \item Transliterate non-Latin scripts if possible; retain diacritics accurately.
\end{itemize}
\end{frame}

\begin{frame}{General Principles (APA / biblatex)}
\begin{itemize}
  \item Use the author’s name in the script of publication (APA allows Latin transliteration).
  \item Give publication data exactly as printed (year, publisher).  APA 7
    dropped the publisher's location.
  \item If the reader is unlikely to understand the title:
    \begin{itemize}
      \item Add a translation in square brackets, outside the italics:
        \emph{Válka s Mloky} [War with the Newts].
    \end{itemize}
  \item Don’t invent English titles — translate accurately but informally.
  \item If the work has an official English edition, you may cite that instead or alongside.
\end{itemize}
\end{frame}

\begin{frame}[fragile]{Example: Czech Source (Čapek, 1936)}
\small
\begin{verbatim}
@book{capek1936,
  author    = {Čapek, Karel},
  year      = {1936},
  title     = {Válka s Mloky [War with the Newts]},
  publisher = {Fr. Borový},
  langid    = {czech}
}
\end{verbatim}

\textbf{Text citation examples:}
\begin{itemize}
  \item \citet{capek1936} satirizes industrial modernity through the figure of the salamander.
  \item The allegory of human exploitation appears early in the narrative \citep{capek1936}.
\end{itemize}

\textbf{References entry (APA style):}
\begin{quote}
Čapek, K. (1936). \emph{Válka s Mloky} [War with the Newts]. Fr. Borový.
\end{quote}
{\footnotesize No place of publication: APA 7 dropped it.  Put the
translation in the \texttt{title} field --- biblatex has a
\texttt{titleaddon} field, but this style brackets it for you, so typing your
own brackets there gives you two sets.}
\end{frame}

\begin{frame}{Example: Japanese Source (\cjk{芥川龍之介}, 1918)}
\small
{\ttfamily
@article\{akutagawa1918,\\
\hspace*{2em}author = \{\cjk{芥川龍之介}\},
\hspace*{2em}sortname = \{Akutagawa, Ryūnosuke\}, \\
\hspace*{2em}year = \{1918\},\\
\hspace*{2em}title = \{\cjk{蜘蛛の糸} Kumo no ito [The Spider's Thread]\},\\
\hspace*{2em}journaltitle = \{\cjk{赤い鳥} Akai Tori\},\\
\hspace*{2em}volume = \{1\}, number = \{1\},\\
\hspace*{2em}langid = \{japanese\}\}\\
}
\vspace{1ex}

\textbf{Text citation examples:}
\begin{itemize}
  \item \citet{akutagawa1918} retells a Buddhist parable of redemption and failure.
  \item Compassion and egoism intertwine in \cjk{「蜘蛛の糸」} \citep{akutagawa1918}.
\end{itemize}

\textbf{References entry (APA 7):}
\begin{quote}
\cjk{芥川龍之介} [Akutagawa, R.]. (1918). \cjk{蜘蛛の糸} \emph{Kumo no ito}
[The spider's thread]. \emph{\cjk{赤い鳥} Akai Tori}, \emph{1}(1).
\end{quote}


\end{frame}



\begin{frame}{When to Translate or Transliterate}
\begin{itemize}
  \item If the audience reads the language → keep original title only.
  \item If not → add translation in brackets after the original title.
  \item Transliteration (romaji, pinyin, etc.) helps with alphabetization and search.
  \item Example (Japanese romanization):
    \begin{quote}\small
    Akutagawa, Ryūnosuke. (1918). \emph{Kumo no ito} [The spider's thread].
    \end{quote}
  \item Always apply one consistent pattern for all non-English items.
\end{itemize}
\end{frame}

% \begin{frame}{In-Text Citation Behavior}
% \begin{itemize}
%   \item \textbf{biblatex} handles Unicode authors automatically.
%   \item In multilingual documents, citations appear correctly:
%     \begin{itemize}
%       \item \citet{capek1936} → Čapek (1936)
%       \item \citet{akutagawa1918} → 芥川龍之介 (1918)
%     \end{itemize}
%   \item Sorting by name or year still works with non-Latin scripts.
%   \item To force alphabetization by Latin transcription, add:  
%     \verb|sortname = {Capek, Karel}| or \verb|sortname = {Akutagawa, Ryunosuke}|.
% \end{itemize}
% \end{frame}

% \begin{frame}{Checklist for Citing Non-English Sources}
% \begin{itemize}
%   \item Verify:
%     \begin{itemize}
%       \item Accurate author spelling and diacritics.
%       \item Year, publisher, and city of publication.
%       \item Correct script and optional translation.
%     \end{itemize}
%   \item Decide: original vs.\ translated title (or both).
% %  \item Add \verb|sortname| for consistent alphabetical order.
% %  \item Test output with \verb|\printbibliography|.
%   \item Keep consistency across all non-English entries.
% \end{itemize}
% \end{frame}

% \begin{frame}{Key Takeaways}
% \begin{itemize}
%   \item Cite foreign-language works with the same rigor as English sources.
%   \item Use original titles + bracketed translations where helpful.
%   \item Unicode and modern \textbf{biblatex} make multilingual citation smooth.
%   \item Careful formatting demonstrates both linguistic and scholarly competence.
%   \item Respect each language’s orthography while following APA consistency.
% \end{itemize}
% \end{frame}


% =====================================================================


\begin{frame}[allowframebreaks]{Acknowledgements}

  \begin{itemize}
  \item The first sections were based on A Short Handbook for Writing Essays in the Humanities and Social Sciences \citep{writinghandbook}
  \item I also consulted \citet{Reid:2024} and \citet{Krause:2007}
  \item \textcite{openai-chatgpt-2025-03-14} was used to format the references, and generate a first draft of the slides with a prompt like
    \begin{flushleft}\ttfamily
       Please make me some slides, based on Chapter 1: https://mlpp.pressbooks.pub/writinghandbook/chapter/chapter-1/ \\

Make them in LaTeX, using luatex and biber, please make around 15 slides with at least 5 bullet points, use sub-lists where appropriate. \\
    \end{flushleft}
  \aicheck
  \end{itemize}
  
\end{frame}


% ====================================================
\begin{frame}[allowframebreaks]
\frametitle{References}
\printbibliography[heading=none]
\end{frame}


% ====================================================
% Appendix: material covered elsewhere, or in more detail
% than the session has time for.  Here for the handout.
\AtBeginSection[]{}
\appendix

\begin{frame}{Appendix}
These frames are here for reference: they repeat material from
another session, or go further than we had time for.
\end{frame}

\begin{frame}{Evaluating Sources — Why It Matters}
\begin{itemize}
  \item Academic writing depends on \textbf{credible evidence}.
  \item Poor sources weaken even the best reasoning.
  \item Evaluating sources ensures:
    \begin{itemize}
      \item Accuracy and reliability
      \item Awareness of bias and limits
      \item Relevance to your argument
    \end{itemize}
  \item Evaluation is a \textbf{critical thinking skill}, not a checklist exercise.
  \end{itemize}

  \vfill
 Adapted from \href{https://wac.colostate.edu/resources/writing/guides/evaluating/}{Evaluating Sources} (WAC Clearinghouse, Colorado State University).  
\end{frame}

\begin{frame}{Purpose and Audience}
\begin{itemize}
  \item Ask: Why was this text created? For whom?
  \item Purposes may include:
    \begin{itemize}
      \item Informing or teaching
      \item Persuading or advocating
      \item Selling or entertaining
    \end{itemize}
  \item Identify intended audience: scholars, professionals, or the general public.
  \item Match the source’s aim with your own research goal.
\end{itemize}
\end{frame}

\begin{frame}{Author and Authority}
\begin{itemize}
  \item Who is the author, and what makes them credible?
  \item Check:
    \begin{itemize}
      \item Education and institutional affiliation
      \item Prior publications and expertise
      \item Reputation in the field
    \end{itemize}
  \item Anonymous or uncredentialed authors demand extra scrutiny.
  \item Authority may also stem from collective or institutional authorship.
\end{itemize}
\end{frame}

\begin{frame}{Publisher and Venue}
\begin{itemize}
  \item Who publishes or hosts the source?
  \item University presses and peer-reviewed journals usually signal quality control.
  \item For websites, assess the domain and hosting organization.
  \item Recognize potential \textbf{institutional bias} in think tanks, corporations, or advocacy groups.
  \item Prefer sources with transparent editorial oversight.
\end{itemize}
\end{frame}

\begin{frame}{Currency and Timeliness}
\begin{itemize}
  \item Consider when the source was written or updated.
  \item In fast-moving fields, information may age quickly.
  \item For historical or theoretical work, older sources may remain foundational.
  \item Look for revision dates, update logs, or newer editions.
  \item Always relate publication date to your topic’s context.
  \item Has the paper been \textbf{retracted}?
    \begin{itemize}
    \item Search in the official \href{https://retractiondatabase.org}{Retraction Watch Database}
    \item Run a search by title, DOI, author, or journal.
    \item This is the most comprehensive open database of retracted papers, maintained by Retraction Watch and made openly available through Crossref.
    \item It also notes the reason for retraction (e.g., plagiarism, data falsification, honest error).
    \end{itemize}
\end{itemize}
\end{frame}

\begin{frame}{Evidence and Support}
\begin{itemize}
  \item Reliable sources \textbf{show their work}.
  \item Ask:
    \begin{itemize}
      \item What kinds of evidence are used? (data, quotations, examples)
      \item Are sources cited and traceable?
      \item Is reasoning logical and transparent?
    \end{itemize}
  \item Unsupported claims or missing citations signal weakness.
  \item Cross-check evidence against other reputable works.
    \begin{itemize}
    \item Multiple sources are more reliable
    \end{itemize}
\end{itemize}
\end{frame}

\begin{frame}{Bias and Objectivity}
\begin{itemize}
  \item No source is completely neutral.
  \item Look for:
    \begin{itemize}
      \item Loaded language or emotional tone
      \item Selective omission of evidence
      \item Conflicts of interest or funding ties
    \end{itemize}
  \item Identify perspective; judge how it shapes interpretation.
  \item Acknowledge bias rather than ignoring it.
\end{itemize}
\end{frame}

\begin{frame}{Balance and Completeness}
\begin{itemize}
  \item Does the source present multiple viewpoints fairly?
  \item Recognize one-sided or partial presentations.
    \begin{itemize}
    \item If they cite mostly themselves it is a bad sign (what counts as
      ``too much'' varies a lot by field)
    \item If they only cite their colleagues it is a bad sign
    \end{itemize}
  \item Check whether evidence contradicting the claim is addressed.
  \item Balanced sources strengthen your own credibility when cited.
  \item Even biased sources can be useful if analyzed critically.
\end{itemize}
\end{frame}

\begin{frame}{Relevance to Your Project}
\begin{itemize}
  \item Determine how the source connects to your research question.
  \item Directly relevant sources:
    \begin{itemize}
      \item Support or challenge your thesis
      \item Provide key evidence or theory
    \end{itemize}
  \item Peripheral sources may supply context or background.
  \item Avoid citing tangential material to inflate your bibliography.
\end{itemize}
\end{frame}

\begin{frame}{Primary vs.\ Secondary Sources}
\begin{itemize}
  \item \textbf{Primary:} original materials (texts, data, interviews, artifacts).
  \item \textbf{Secondary:} analysis, interpretation, commentary.
  \item  \textbf{Tertiary:} index or textual consolidation of  primary and secondary sources
  \item Choose according to purpose:
    \begin{itemize}
      \item Primary for direct evidence
      \item Secondary for framing and critique
      \item Tertiary for an overview
    \end{itemize}
  \item Distinguish between firsthand and filtered perspectives.
\end{itemize}
\end{frame}

\begin{frame}{Scholarly vs.\ Popular Sources}
\begin{itemize}
  \item \textbf{Scholarly:} peer-reviewed, technical, detailed references.
  \item \textbf{Popular:} general readership, journalistic style.
  \item Use scholarly works for evidence, popular for public context.
  \item Be cautious: some “grey literature” mixes the two.
  \item Evaluate tone, citations, and rigor to tell them apart.
\end{itemize}
\end{frame}

\begin{frame}{Checklist for Evaluating a Source}
\begin{itemize}
  \item Purpose and audience clearly stated?
  \item Author’s credentials and affiliations verifiable?
  \item Publisher or host credible and transparent?
  \item Evidence traceable and balanced?
  \item Date current enough for the topic?
  \item Bias recognized and context considered?
  \item Source relevant to your own argument?
\end{itemize}
\end{frame}

\begin{frame}{Integrating Sources}
\begin{itemize}
  \item Use quotation, paraphrase, and summary effectively.
  \item Cite sources to:
    \begin{itemize}
      \item Credit others’ ideas
      \item Strengthen your credibility
      \item Help readers locate materials
    \end{itemize}
  \item Follow the style your community uses --- APA 7 in this course.
  \item Blend sources seamlessly with your own analysis.
\end{itemize}
\end{frame}


\end{document}
%%% Local Variables:
%%% coding: utf-8
%%% mode: latex
%%% TeX-engine: luatex
%%% End:
